% Préambule partagé pour les fiches de révision et exercices générés à partir des cours.
% Ce fichier est copié dans chaque dossier de cours (~/Cours/<id>/preambule.tex) puis
% importé avec % Préambule partagé pour les fiches de révision et exercices générés à partir des cours.
% Ce fichier est copié dans chaque dossier de cours (~/Cours/<id>/preambule.tex) puis
% importé avec % Préambule partagé pour les fiches de révision et exercices générés à partir des cours.
% Ce fichier est copié dans chaque dossier de cours (~/Cours/<id>/preambule.tex) puis
% importé avec % Préambule partagé pour les fiches de révision et exercices générés à partir des cours.
% Ce fichier est copié dans chaque dossier de cours (~/Cours/<id>/preambule.tex) puis
% importé avec \input{preambule.tex} pour que chaque dossier reste autonome et compilable
% indépendamment avec pdflatex.

\usepackage[margin=2.2cm]{geometry}
\usepackage[utf8]{inputenc}
\usepackage[T1]{fontenc}

% mathématiques
\usepackage{amsmath}
\usepackage{amsfonts}
\usepackage{amssymb}
\usepackage{mathtools}
\usepackage{bm}
\def\vec#1{{\ensuremath{\bm{{#1}}}}}

\DeclareMathOperator*{\argmax}{arg\,max}
\DeclareMathOperator*{\argmin}{arg\,min}
\newcommand{\R}{\mathbb{R}}
\newcommand{\E}{\mathbb{E}}
\newcommand{\PP}{\mathbb{P}}
\newcommand{\N}{\mathbb{N}}

% liens et mise en page
\usepackage[colorlinks=true, linkcolor=blue, urlcolor=blue]{hyperref}
\usepackage{graphicx}
\usepackage{enumitem}
\usepackage{parskip}
\usepackage{xcolor}
\usepackage{fancyhdr}

\pagestyle{fancy}
\fancyhf{}
\fancyhead[L]{\small\leftmark}
\fancyfoot[C]{\thepage}
\renewcommand{\headrulewidth}{0.3pt}

% boîtes pour concepts clés / définitions
\usepackage[most]{tcolorbox}
\newtcolorbox{definition}[1][]{
  colback=blue!4, colframe=blue!35!black, boxrule=0.6pt, arc=2pt,
  left=8pt, right=8pt, top=6pt, bottom=6pt, #1
}
\newtcolorbox{correction}[1][]{
  colback=green!4, colframe=green!30!black, boxrule=0.6pt, arc=2pt,
  left=8pt, right=8pt, top=6pt, bottom=6pt, #1
}
 pour que chaque dossier reste autonome et compilable
% indépendamment avec pdflatex.

\usepackage[margin=2.2cm]{geometry}
\usepackage[utf8]{inputenc}
\usepackage[T1]{fontenc}

% mathématiques
\usepackage{amsmath}
\usepackage{amsfonts}
\usepackage{amssymb}
\usepackage{mathtools}
\usepackage{bm}
\def\vec#1{{\ensuremath{\bm{{#1}}}}}

\DeclareMathOperator*{\argmax}{arg\,max}
\DeclareMathOperator*{\argmin}{arg\,min}
\newcommand{\R}{\mathbb{R}}
\newcommand{\E}{\mathbb{E}}
\newcommand{\PP}{\mathbb{P}}
\newcommand{\N}{\mathbb{N}}

% liens et mise en page
\usepackage[colorlinks=true, linkcolor=blue, urlcolor=blue]{hyperref}
\usepackage{graphicx}
\usepackage{enumitem}
\usepackage{parskip}
\usepackage{xcolor}
\usepackage{fancyhdr}

\pagestyle{fancy}
\fancyhf{}
\fancyhead[L]{\small\leftmark}
\fancyfoot[C]{\thepage}
\renewcommand{\headrulewidth}{0.3pt}

% boîtes pour concepts clés / définitions
\usepackage[most]{tcolorbox}
\newtcolorbox{definition}[1][]{
  colback=blue!4, colframe=blue!35!black, boxrule=0.6pt, arc=2pt,
  left=8pt, right=8pt, top=6pt, bottom=6pt, #1
}
\newtcolorbox{correction}[1][]{
  colback=green!4, colframe=green!30!black, boxrule=0.6pt, arc=2pt,
  left=8pt, right=8pt, top=6pt, bottom=6pt, #1
}
 pour que chaque dossier reste autonome et compilable
% indépendamment avec pdflatex.

\usepackage[margin=2.2cm]{geometry}
\usepackage[utf8]{inputenc}
\usepackage[T1]{fontenc}

% mathématiques
\usepackage{amsmath}
\usepackage{amsfonts}
\usepackage{amssymb}
\usepackage{mathtools}
\usepackage{bm}
\def\vec#1{{\ensuremath{\bm{{#1}}}}}

\DeclareMathOperator*{\argmax}{arg\,max}
\DeclareMathOperator*{\argmin}{arg\,min}
\newcommand{\R}{\mathbb{R}}
\newcommand{\E}{\mathbb{E}}
\newcommand{\PP}{\mathbb{P}}
\newcommand{\N}{\mathbb{N}}

% liens et mise en page
\usepackage[colorlinks=true, linkcolor=blue, urlcolor=blue]{hyperref}
\usepackage{graphicx}
\usepackage{enumitem}
\usepackage{parskip}
\usepackage{xcolor}
\usepackage{fancyhdr}

\pagestyle{fancy}
\fancyhf{}
\fancyhead[L]{\small\leftmark}
\fancyfoot[C]{\thepage}
\renewcommand{\headrulewidth}{0.3pt}

% boîtes pour concepts clés / définitions
\usepackage[most]{tcolorbox}
\newtcolorbox{definition}[1][]{
  colback=blue!4, colframe=blue!35!black, boxrule=0.6pt, arc=2pt,
  left=8pt, right=8pt, top=6pt, bottom=6pt, #1
}
\newtcolorbox{correction}[1][]{
  colback=green!4, colframe=green!30!black, boxrule=0.6pt, arc=2pt,
  left=8pt, right=8pt, top=6pt, bottom=6pt, #1
}
 pour que chaque dossier reste autonome et compilable
% indépendamment avec pdflatex.

\usepackage[margin=2.2cm]{geometry}
\usepackage[utf8]{inputenc}
\usepackage[T1]{fontenc}

% mathématiques
\usepackage{amsmath}
\usepackage{amsfonts}
\usepackage{amssymb}
\usepackage{mathtools}
\usepackage{bm}
\def\vec#1{{\ensuremath{\bm{{#1}}}}}

\DeclareMathOperator*{\argmax}{arg\,max}
\DeclareMathOperator*{\argmin}{arg\,min}
\newcommand{\R}{\mathbb{R}}
\newcommand{\E}{\mathbb{E}}
\newcommand{\PP}{\mathbb{P}}
\newcommand{\N}{\mathbb{N}}

% liens et mise en page
\usepackage[colorlinks=true, linkcolor=blue, urlcolor=blue]{hyperref}
\usepackage{graphicx}
\usepackage{enumitem}
\usepackage{parskip}
\usepackage{xcolor}
\usepackage{fancyhdr}

\pagestyle{fancy}
\fancyhf{}
\fancyhead[L]{\small\leftmark}
\fancyfoot[C]{\thepage}
\renewcommand{\headrulewidth}{0.3pt}

% boîtes pour concepts clés / définitions
\usepackage[most]{tcolorbox}
\newtcolorbox{definition}[1][]{
  colback=blue!4, colframe=blue!35!black, boxrule=0.6pt, arc=2pt,
  left=8pt, right=8pt, top=6pt, bottom=6pt, #1
}
\newtcolorbox{correction}[1][]{
  colback=green!4, colframe=green!30!black, boxrule=0.6pt, arc=2pt,
  left=8pt, right=8pt, top=6pt, bottom=6pt, #1
}
